\documentclass[11pt]{article}
\usepackage{amsmath,amssymb,amsthm,hyperref}
\newtheorem{theorem}{Theorem}
\newtheorem{proposition}{Proposition}
\newtheorem{corollary}{Corollary}
\newcommand{\norm}[1]{\left\lVert #1\right\rVert}
\newcommand{\Real}{\operatorname{Re}}
\newcommand{\sgn}{\operatorname{sgn}}
\title{Real signings from Hermitian orientations through odd-cycle maps}
\author{Independent signing session}
\date{15 November 2026}
\begin{document}
\maketitle

\noindent\textbf{Status.} The proofs below establish a restricted signing
theorem and a limitation of its mechanism. They do not establish a
$2+o_g(1)$ theorem for all high-girth regular graphs. Historical novelty
has not been established. No general Bilu--Linial conjecture is assumed.

For a finite simple graph $G$, let $\mu_G$ be its matching polynomial and
let $r_G$ denote its largest root. Greaves--Mohar--O~\cite[Theorem 1.1,
Corollary 1.2]{GMO} prove that there is a real skew-symmetric matrix $S$
with entries $S_{uv}\in\{\pm1\}$ on edges, zero elsewhere, such that
$H=iS$ is Hermitian and $\norm{H}\le r_G$. Their interlacing family is
an orientation family, and its spectra are symmetric about zero.
This existing theorem supplies all interlacing used below.

\begin{theorem}[Odd-cycle map conversion]\label{thm:cycle}
Suppose $L\ge3$ is odd and $G$ admits a graph homomorphism to $C_L$.
Then $G$ has a real edge signing $\sigma$ satisfying
\[
 \norm{A_\sigma}\le \sec\!\left(\frac{\pi}{2L}\right)r_G.
 \tag{1}
\]
In particular, if $G$ has maximum degree at most $d\ge2$, then
\[
 \norm{A_\sigma}\le
 2\sec\!\left(\frac{\pi}{2L}\right)\sqrt{d-1}.
 \tag{2}
\]
The constant in (1) is best possible over graphs admitting a map to $C_L$.
\end{theorem}
\begin{proof}
Write $L=2m+1$, $\phi=\pi m/L$, and
$c=\sin\phi=\cos(\pi/(2L))$. Let $a:V(G)\to\{0,\ldots,L-1\}$
be the cycle map, so adjacent values differ by $\pm1$ modulo $L$.
Set $D_{vv}=e^{i\phi a(v)}$ and choose $H=iS$ as above. On an edge,
\begin{align*}
 K_{uv}:=(D^*HD)_{uv}&=iS_{uv}e^{i\phi(a(v)-a(u))},\\
 \Real K_{uv}&=-S_{uv}\sin\bigl(\phi(a(v)-a(u))\bigr).
\end{align*}
The integer difference $a(v)-a(u)$ is $\pm1$ or $\pm(L-1)$.
Since $\phi L=\pi m$, every displayed sine has absolute value $c$.
Thus $A_\sigma=c^{-1}\Real K$ is a real symmetric signed adjacency
matrix on exactly $G$. Entrywise real part satisfies
\[
 \norm{\Real K}
 =\norm{(K+\overline K)/2}
 \le\tfrac12(\norm K+\norm{\overline K})=\norm H\le r_G.
\]
This proves (1); the usual matching-root bound
$r_G\le2\sqrt{d-1}$ gives (2).

For sharpness of (1), take $G=C_L$. Every signing is switching equivalent
to a cycle with flux $+1$ or $-1$. The $+1$ class has eigenvalue $2$;
the $-1$ class has eigenvalue $-2$ because $L$ is odd. Both have norm $2$.
Also $\mu_{C_L}(x)=2T_L(x/2)$, so
$r_{C_L}=2\cos(\pi/(2L))$. Equality therefore holds in (1).
\end{proof}

\begin{corollary}
Every $3$-colorable graph of maximum degree $d\ge2$ has a signing of
norm at most $(4/\sqrt3)\sqrt{d-1}$. More generally, for any family of
$d$-regular graphs supplied with maps to $C_{L_j}$, where odd $L_j\to\infty$,
\[
 \min_\sigma\norm{A_\sigma}/\sqrt{d-1}
 \le 2\sec\!\left(\frac{\pi}{2L_j}\right)
 =2+\frac{\pi^2}{4L_j^2}+O(L_j^{-4}).
\]
The statement remains uniform if $d$ varies with $j$.
\end{corollary}

This is a different selection mechanism from recovery from the
$\chi[A,-A]$ positive measure: choose a small-norm imaginary Hermitian
orientation first and take a scaled real part after a fixed vertex gauge.
For example, the coefficient for a map to $C_5$ is
$2\sec(\pi/10)=2.102924\ldots<3/\sqrt2$.
No girth assumption is required for this restricted conclusion.

\begin{proposition}[Exact applicability of constant-modulus phases]
Let $G$ be finite, connected, and contain an edge. Let $c_*(G)$ be the
supremum of $c\in[0,1]$ for which phases satisfy
$|\sin(\theta_v-\theta_u)|=c$ on every edge. If $G$ is bipartite,
$c_*(G)=1$. If $G$ is nonbipartite and $3$-colorable, define
\[
 L(G)=\max\{L\ge3:\ L\text{ odd and }G\to C_L\}.
\]
Then $L(G)$ exists and $c_*(G)=\cos(\pi/(2L(G)))$.
If $\chi(G)\ge4$, then $c_*(G)=0$.
\end{proposition}
\begin{proof}
For any $c>0$ put $\alpha=\arcsin c\in(0,\pi/2]$. Modulo $\pi$,
every edge difference is $\pm\alpha$. Relative to a root, all vertex
phases therefore belong to the subgroup generated by $\alpha$ in
$\mathbb R/(\pi\mathbb Z)$. If $\alpha/\pi$ is irrational, this
subgroup is infinite cyclic, and integer labels with adjacent difference
$\pm1$ give a bipartition. If $\alpha/\pi=a/b$ in lowest terms, its
order is $b$, and the same labels modulo $b$ give a cycle homomorphism.
Even $b$ again gives a bipartition. For a nonbipartite graph, necessarily
$b\ge3$ is odd and $G\to C_b$.

Fix any odd cycle of length $\ell$ in $G$. Its image in $C_b$ is an
$\ell$-step closed walk with integer displacement an odd multiple of $b$.
Thus $b\le\ell$, proving finiteness of the set of possible odd $b$.
For each such denominator,
$|\sin(\pi a/b)|\le\cos(\pi/(2b))\le\cos(\pi/(2L(G)))$.
The phase assignment in Theorem~\ref{thm:cycle} attains equality.
Every $3$-colorable graph maps to $C_3$, whereas a graph mapping to any
odd cycle is $3$-colorable. This proves all nonbipartite assertions.
For a bipartite graph, phases $0,\pi/2$ on its two parts attain $c=1$.
\end{proof}

Consequently, optimizing all scalar vertex gauges and a single constant
rescaling gives exactly this cycle-map mechanism. It does not secretly
extend to additional nonbipartite graphs.

\begin{proposition}[Obstruction to scalar phase conversion]\label{prop:obstruction}
Let $G$ be a finite $d$-regular graph with unsigned adjacency $A$ and
least eigenvalue $\lambda_{\min}$. Suppose vertex phases $\theta_v$ and
$c\in[0,1]$ obey
\[
 |\sin(\theta_v-\theta_u)|\ge c\quad\text{on every edge.}
 \tag{3}
\]
Then
\[
 c^2\le \frac{1-\lambda_{\min}/d}{2}.
 \tag{4}
\]
In particular, if a single phase gauge is to turn every edge of $iS$
into a real part of constant absolute value $c$, a necessary condition
for a conversion loss $1/c=1+o(1)$ is
$\lambda_{\min}(A)/d=-1+o(1)$.
\end{proposition}
\begin{proof}
Put $z_v=e^{2i\theta_v}$. Then $\norm z^2=|V(G)|=n$, and
\[
 z^*Az=2\sum_{uv\in E(G)}\cos(2(\theta_v-\theta_u))
 \le nd(1-2c^2).
\]
The variational bound $z^*Az\ge n\lambda_{\min}$ proves (4).
The final assertion follows since $\lambda_{\min}\ge-d$.
\end{proof}

For example, on any $d$-regular graph with
$\lambda_{\min}\ge-2\sqrt{d-1}-\varepsilon$, every constant-modulus
phase conversion has loss at least
\[
 \frac1c\ge
 \sqrt{\frac{2d}{d+2\sqrt{d-1}+\varepsilon}}.
 \tag{5}
\]
For each fixed $d\ge3$ and $\varepsilon\to0$, this bound is strictly
larger than $1$. As $d\to\infty$ it tends to $\sqrt2$.
Equation (5) limits the conversion factor itself; an orientation could
have a smaller norm than the generic matching bound, so it is not a
lower bound for the true best signing norm.

There is an additional elementary obstruction: $G\to C_L$ entails
$\chi(G)\le3$. High girth alone does not force $3$-colorability. Even
within bounded degree, the spectral condition above shows that a long
odd-cycle map requires the unsigned graph to be spectrally close to
bipartite. In fact (4) with $c=\cos(\pi/(2L))$ gives
\[
 \lambda_{\min}(A)\le-d\cos(\pi/L).
 \tag{6}
\]
Thus a large-girth result cannot be deduced by silently replacing
girth with the odd-cycle map parameter.

\paragraph{Further justified extension.}
If removing an edge set $F$ yields $G_0\to C_L$ and the adjacency matrix
of $F$ has maximum degree $r$, sign $G_0$ by Theorem~\ref{thm:cycle}
and sign $F$ arbitrarily. The elementary row-sum bound gives
\[
 \norm{A_\sigma(G)}\le\sec(\pi/(2L))r_{G_0}+r.
\]
This only helps when $r/\sqrt{d-1}$ is small; the note does not claim
that high girth provides such an edge set.

\paragraph{Exact unresolved step.}
For arbitrary high-girth graphs, the required cycle map or an equally
strong norm-preserving real conversion is absent. Proposition
\ref{prop:obstruction} prevents fixing that gap by optimizing a single
vertex phase gauge and constant entrywise scaling. Neither the theorem
nor its obstruction addresses all matrix-valued gauges, nonconstant
Schur multipliers, or simultaneous constrained root selection.

\paragraph{Primary-source comparison and priority scope.}
The existence of the Hermitian orientation and its interlacing proof are
entirely due to Greaves--Mohar--O~\cite{GMO}; the real-signing conversion
above is proved here as an immediate additional construction. Searches
on 15 November 2026 for graph signing with odd-cycle homomorphisms,
Hermitian phase conversion, and $3$-colorable graph signings found no
matching statement. This is insufficient to establish historical novelty.
Marcus--Spielman--Srivastava~\cite{MSS} supply one-sided real-signing
control and the matching bound. Lin--Zhou~\cite{LZ} use a distinct
mixed-root recovery argument, giving the general coefficient
$(3+\sqrt5)/2$ and short-cycle improvements. The present statement
requires the extra cycle-map hypothesis, so it is not a general
improvement of their theorem. Xu~\cite{Xu} disproves exact universal
real signing at coefficient $2$; this note does not conflict with that
result. The supplied friend's $3/\sqrt2+O(1/g)$ announcement and
degree-$32$ hardness announcement were not used as theorems.

There is also a preexisting near-Ramanujan theorem with a different
hypothesis: Mohanty--O'Donnell--Paredes~\cite[Theorem 3.1]{MOP} prove,
for $d\le\operatorname{polylog}n$ and bicycle-free radius
$r\gg(\log\log n)^2$, that random signs have adjacency norm at most
$2\sqrt{d-1}(1+O((\log\log n)^4/r^2))$ with probability at least
$1-n^{-100}$. This has an explicit dependence on graph order through
its radius hypothesis. Our restricted cycle-map theorem is independent
of order but does not remove their hypothesis for arbitrary graphs.

\begin{thebibliography}{9}
\bibitem{GMO} G. Greaves, B. Mohar, and S. O,
\emph{Interlacing families and the Hermitian spectral norm of digraphs},
\url{https://arxiv.org/html/1602.06274}.
Theorem 1.1, Corollaries 1.2--1.3, and Theorem 3.4 read directly.
\bibitem{MSS} A. Marcus, D. Spielman, and N. Srivastava,
\emph{Interlacing families I: Bipartite Ramanujan graphs of all degrees},
\url{https://arxiv.org/abs/1304.4132}.
\bibitem{LZ} F. Lin and H. Zhou,
\emph{Improved Bounds for the Bilu--Linial Conjecture via Spectral Recovery
from Mixed Determinantal Polynomials}, version 2,
\url{https://arxiv.org/html/2609.15715v2}.
\bibitem{Xu} Z. Xu,
\emph{A $3$-regular counterexample to the Bilu--Linial signing conjecture},
version 4, \url{https://arxiv.org/html/2609.15591v4}.
\bibitem{MOP} S. Mohanty, R. O'Donnell, and P. Paredes,
\emph{Explicit near-Ramanujan graphs of every degree},
\url{https://arxiv.org/html/1909.06988}, Theorem 3.1.
\end{thebibliography}
\end{document}
